\providecommand{\UseMathForPositioningText}{}
\documentclass[11pt,a4paper]{article}

\usepackage{geometry}
\geometry{top=2.5cm,bottom=2.5cm,left=2.5cm,right=2.5cm}
\usepackage{graphicx}
\usepackage{setspace}
\usepackage{tikz}
\usetikzlibrary{calc}
\usetikzlibrary{arrows.meta,positioning}
\usepackage{eso-pic}
\usepackage{titlesec}
\usepackage{float}
\usepackage{amsmath}
\usepackage{amssymb}
\usepackage{booktabs}
\usepackage{array}
\usepackage{caption}
\usepackage{tocloft}
\usepackage{xcolor}
\usepackage{hyperref}
\usepackage{listings}
\usepackage{enumitem}
\usepackage{longtable}
\usepackage{tabularx}
\usepackage{pifont}
\usepackage{xepersian}

\IfFontExistsTF{XB Niloofar}
  {\settextfont{XB Niloofar}}
  {\settextfont{Noto Naskh Arabic}}

\IfFontExistsTF{DejaVu Sans}
  {\setlatintextfont{DejaVu Sans}}
  {}

\IfFontExistsTF{DejaVu Sans Mono}
  {\setmonofont{DejaVu Sans Mono}}
  {}

\titleformat{\section}{\fontsize{18pt}{22pt}\bfseries}{\thesection}{1em}{}
\titleformat{\subsection}{\fontsize{15pt}{19pt}\bfseries}{\thesubsection}{1em}{}
\titleformat{\subsubsection}{\fontsize{13pt}{17pt}\bfseries}{\thesubsubsection}{1em}{}
\renewcommand{\cftsecleader}{\cftdotfill{\cftdotsep}}

\setstretch{1.25}
\setlength{\parindent}{1.2cm}
\setlength{\parskip}{0.35em}
\setlength{\emergencystretch}{2em}
\setlength{\footnotesep}{0.5cm}

\newcommand{\enfoot}[1]{\LTRfootnote{#1}}
\newcommand{\code}[1]{\lr{\texttt{\detokenize{#1}}}}
\newcommand{\filename}[1]{\lr{\texttt{\detokenize{#1}}}}
\newcommand{\passmark}{\textcolor{black}{\ding{51}}}
\renewcommand{\labelitemi}{\(\bullet\)}
\renewcommand{\lstlistingname}{کد}

\lstdefinestyle{verilogstyle}{
  language=Verilog,
  basicstyle=\ttfamily\small,
  numbers=left,
  numberstyle=\tiny,
  stepnumber=1,
  numbersep=8pt,
  frame=single,
  rulecolor=\color{black!45},
  backgroundcolor=\color{black!3},
  keywordstyle=\color{blue!55!black}\bfseries,
  commentstyle=\color{green!35!black},
  stringstyle=\color{red!55!black},
  showstringspaces=false,
  breaklines=true,
  breakatwhitespace=false,
  tabsize=4,
  columns=fullflexible,
  keepspaces=true,
  captionpos=b
}
\lstset{style=verilogstyle}

\newcommand{\reportimage}[2][0.96\textwidth]{%
  \IfFileExists{#2}%
    {\includegraphics[width=#1]{#2}}%
    {\fbox{\parbox[c][5.2cm][c]{0.88\textwidth}{\centering\textbf{محل قرارگیری تصویر}\\[0.35cm]\lr{\texttt{#2}}}}}%
}

\begin{document}
\AddToShipoutPictureBG{
  \begin{tikzpicture}[remember picture,overlay]
    \node[
      draw,
      line width=2.5pt,
      minimum width=\paperwidth-2cm,
      minimum height=\paperheight-2cm,
      inner sep=0pt
    ] at (current page.center) {};
    \node[
      draw,
      line width=1pt,
      minimum width=\paperwidth-2.3cm,
      minimum height=\paperheight-2.3cm,
      inner sep=0pt
    ] at (current page.center) {};
  \end{tikzpicture}
}

\begin{titlepage}
  \centering
  \vspace*{0.1cm}

  \IfFileExists{images/logo.png}
    {
\includegraphics[width=2.9cm]{images/logo.png}}
    {\fbox{\parbox[c][2.6cm][c]{2.6cm}{\centering\lr{images/logo.png}}}}
  \\[0.35cm]
  {\fontsize{13pt}{17pt}\selectfont دانشگاه صنعتی شریف}\\[0.1cm]
  {\fontsize{13pt}{17pt}\selectfont دانشکده مهندسی کامپیوتر}

  \vspace{0.75cm}
  {\fontsize{15pt}{19pt}\bfseries گزارش کار آزمایش هفتم}

  \vspace{0.7cm}
  {\fontsize{11.5pt}{15pt}\selectfont عنوان:}\\[0.3cm]
  {\fontsize{19pt}{25pt}\bfseries طراحی و پیاده‌سازی پروتکل \lr{UART}}

  \vspace{0.75cm}
  {\fontsize{11.5pt}{15pt}\selectfont نام درس}\\[0.25cm]
  {\fontsize{13.5pt}{17pt}\bfseries آزمایشگاه طراحی سامانه‌های رقمی}

  \vspace{0.6cm}
  {\fontsize{11.5pt}{15pt}\selectfont تابستان ۱۴۰۵}

  \vspace{0.6cm}
  {\fontsize{11.5pt}{15pt}\selectfont استاد}\\[0.2cm]
  {\fontsize{13pt}{17pt}\bfseries دکتر اجلالی}

  \vspace{0.45cm}
  {\fontsize{11.5pt}{15pt}\selectfont دستیار استاد}\\[0.2cm]
  {\fontsize{13pt}{17pt}\bfseries علی اثنی عشری}

  \vspace{0.55cm}
  {\fontsize{11.5pt}{15pt}\selectfont گروه ۲\par}

  \vspace{0.25cm}
  {\fontsize{13pt}{17pt}\bfseries
  علی سلطانی ـ ۴۰۳۱۰۶۰۹۲\\[0.12cm]
  یگانه کربلایی ـ ۴۰۳۱۰۶۵۲۷\\[0.12cm]
  محمدمهدی مرادی ـ ۴۰۳۱۰۶۶۸۱\par}

  \vfill
\end{titlepage}

\newpage
\tableofcontents
\newpage

\section{هدف آزمایش}

هدف این آزمایش، طراحی و پیاده‌سازی یک واحد ارتباط سریال ناهم‌زمان \lr{UART}\enfoot{Universal Asynchronous Receiver Transmitter} و بررسی عملکرد آن در سطح شبیه‌سازی است. در این طراحی، داده‌های هفت‌بیتی به‌صورت سریال و بر اساس یک قالب قاب مشخّص منتقل می‌شوند. هر قاب شامل یک بیت شروع\enfoot{Start Bit}، یک بیت توازن\enfoot{Parity Bit}، هفت بیت داده و یک بیت توقف\enfoot{Stop Bit} است؛ در نتیجه برای هر داده، در مجموع ۱۰ بیت روی خط سریال قرار می‌گیرد.

هدف اصلی پیاده‌سازی، ساخت یک فرستنده، یک گیرنده و یک پیمانهٔ سطح بالا است که خروجی سریال فرستنده را مستقیماً به ورودی سریال گیرنده متصل کند. به این ساختار، بازگشت مستقیم\enfoot{Loopback} گفته می‌شود. در حالت صحیح، دادهٔ دریافت‌شده باید با دادهٔ ارسال‌شده برابر باشد و سیگنال صحت داده تنها در صورت معتبر بودن قاب فعال شود.

علاوه بر مسیر عادی، فایل آزمون به‌صورت جامع طراحی شده است تا خطاهای توازن و بیت توقف، درخواست داده در زمان مشغول‌بودن فرستنده، بازنشانی ناهم‌زمان در میانهٔ انتقال، پالس شروع کاذب و همچنین تمامی ۱۲۸ مقدار ممکن دادهٔ هفت‌بیتی را بررسی کند. به این ترتیب، درستی مدار تنها با یک نمونهٔ عادی اثبات نمی‌شود و حالت‌های مرزی نیز به‌صورت خودکار آزموده می‌شوند.

\section{شرح مسئله و قالب قاب \lr{UART}}

\subsection{ساختار قاب}

در حالت بی‌کار، خط سریال مقدار ۱ دارد. هنگام آغاز انتقال، فرستنده ابتدا یک صفر را به‌عنوان بیت شروع روی خط قرار می‌دهد. پس از آن بیت توازن، سپس بیت‌های داده از \code{D0} تا \code{D6} و در پایان بیت توقف با مقدار ۱ قرار می‌گیرند. ترتیب کامل قاب به‌صورت زیر است:

\begin{latin}
\begin{center}
\texttt{Idle (1) $\rightarrow$ Start (0) $\rightarrow$ Parity $\rightarrow$ D0 $\rightarrow$ D1 $\rightarrow$ D2 $\rightarrow$ D3 $\rightarrow$ D4 $\rightarrow$ D5 $\rightarrow$ D6 $\rightarrow$ Stop (1)}
\end{center}
\end{latin}

بیت توازن از طریق عملگر \lr{XOR} روی هفت بیت داده محاسبه می‌شود:

\[
P = D_0 \oplus D_1 \oplus D_2 \oplus D_3 \oplus D_4 \oplus D_5 \oplus D_6
\]

برای نمونه، در صورت ارسال دادهٔ

\[
\texttt{send\_data}=7'b1011001
\]

گیرنده باید پس از پایان قاب مقدار \code{rec_data}=\code{1011001} را تولید کند و در صورت صحیح‌بودن بیت توازن و بیت توقف، سیگنال \code{correct_data} را برابر ۱ قرار دهد.

\subsection{زمان‌بندی انتقال}

نرخ استاندارد ذکرشده در صورت آزمایش برابر \lr{115200 bps} است. در این نرخ، زمان یک بیت از رابطهٔ

\[
T_{bit}=\frac{1}{115200}
\]

به‌دست می‌آید. اگر ساعت سامانه ۵۰ مگاهرتز باشد، تعداد چرخه‌های ساعت لازم برای یک بیت تقریباً برابر است با

\[
\frac{50\,000\,000}{115200}\simeq 434.
\]

از آنجا که شبیه‌سازی با هدف مشاهدهٔ سریع‌تر شکل موج انجام شده است، در پیاده‌سازی حاضر زمان هر بیت با پارامتر \code{BIT_TICKS} کنترل می‌شود و مقدار آن برابر ۴ قرار گرفته است. همچنین دورهٔ ساعت در فایل آزمون برابر ۱۰ نانوثانیه انتخاب شده است؛ در نتیجه هر بیت چهار چرخهٔ ساعت، یعنی ۴۰ نانوثانیه، روی خط باقی می‌ماند.

\section{طراحی فرستنده}

\subsection{ساختار ماشین حالت فرستنده}

فرستنده به‌صورت یک ماشین حالت متناهی\enfoot{Finite State Machine} پنج‌حالته پیاده‌سازی شده است. دلیل استفاده از ماشین حالت این است که ارسال یک قاب \lr{UART} را می‌توان به چند مرحلهٔ متوالی و کاملاً مشخّص تقسیم کرد و در هر مرحله مقدار خروجی، شمارندۀ زمان و سیگنال‌های کنترلی را به‌صورت مستقل کنترل نمود. در پیاده‌سازی حاضر این حالت‌ها شامل \code{IDLE}، \code{START}، \code{PARITY}، \code{DATA} و \code{STOP} هستند.

در حالت \code{IDLE}، خط \code{tx} در سطح منطقی ۱ نگه داشته می‌شود و \code{busy} برابر صفر است. در همین حالت، فرستنده منتظر یک پالس \code{new_data} می‌ماند. به‌محض مشاهدهٔ این سیگنال، مقدار هفت‌بیتی \code{send_data} در ثبات داخلی \code{data_reg} لچ می‌شود. این نکته اهمیت زیادی دارد، زیرا از این لحظه به بعد تغییرات روی ورودی خارجی \code{send_data} نباید روی قاب جاری اثر بگذارند. هم‌زمان بیت توازن با عملگر کاهش‌دهندۀ \lr{XOR} روی همان داده محاسبه و در \code{parity_reg} ذخیره می‌شود. سپس \code{busy} فعال شده و خروجی \code{tx} مستقیماً به صفر تغییر می‌کند تا بیت شروع آغاز شود.

در حالت \code{START} مقدار صفر به مدت دقیقاً \code{BIT_TICKS} چرخه روی خط باقی می‌ماند. شمارندۀ \code{tick_count} در هر لبهٔ بالاروندهٔ ساعت یک واحد افزایش پیدا می‌کند و هنگامی که به \code{BIT_TICKS-1} برسد، زمان این بیت کامل شده و شمارنده دوباره صفر می‌شود. در همان لحظه فرستنده وارد حالت \code{PARITY} می‌شود و بیت توازن ذخیره‌شده را روی \code{tx} قرار می‌دهد. به این ترتیب مرز بین بیت شروع و بیت توازن کاملاً به زمان‌بندی ساعت وابسته و ثابت است.

در حالت \code{PARITY}، مقدار \code{parity_reg} برای یک زمان بیت کامل روی خط باقی می‌ماند. پس از پایان این زمان، اندیس بیت \code{bit_idx} صفر می‌شود و نخستین بیت داده یعنی \code{data_reg[0]} روی \code{tx} قرار می‌گیرد. این انتخاب عمداً با ترتیب صورت آزمایش هماهنگ است که بیت‌های داده را از کم‌ارزش‌ترین بیت، یعنی \code{D0}، به‌سوی پرارزش‌ترین بیت، یعنی \code{D6}، ارسال می‌کند.

حالت \code{DATA} بخش تکرارشوندهٔ ماشین حالت است. در هر مرحله، بیت فعلی با عبارت \code{data_reg[bit_idx]} انتخاب می‌شود و تا پایان \code{BIT_TICKS} چرخه روی خط باقی می‌ماند. اگر اندیس هنوز به ۶ نرسیده باشد، پس از پایان زمان بیت، \code{bit_idx} یک واحد افزایش می‌یابد و بیت بعدی انتخاب می‌شود. هنگامی که \code{bit_idx=6} باشد، پس از پایان زمان \code{D6}، فرستنده به حالت \code{STOP} منتقل می‌شود و مقدار \code{tx} برابر یک قرار می‌گیرد.

در حالت \code{STOP} نیز مقدار یک به مدت یک زمان بیت روی خط باقی می‌ماند. این بیت پایان قاب را مشخص می‌کند و تا پایان آن، \code{busy} همچنان یک باقی می‌ماند. تنها پس از کامل‌شدن آخرین زمان بیت است که \code{busy} صفر می‌شود و ماشین حالت به \code{IDLE} بازمی‌گردد. بنابراین از دید مدار بالادست، صفر شدن \code{busy} به‌معنای پایان کامل قاب است و نه صرفاً پایان بخش داده.

یکی دیگر از ویژگی‌های مهم طراحی، بازنشانی ناهم‌زمان است. با فعال‌شدن \code{rstN=0}، حالت به \code{IDLE}، خروجی \code{tx} به ۱، \code{busy} به صفر و شمارنده‌ها به مقدار اولیه بازمی‌گردند. به همین دلیل حتی اگر بازنشانی در وسط یک قاب اعمال شود، فرستنده بدون انتظار برای لبهٔ ساعت بعدی قاب را رها کرده و بلافاصله به وضعیت بی‌کار می‌رود. این ویژگی مستقیماً در آزمون نهم فایل آزمون نیز بررسی شده است.

همچنین، در حالت \code{IDLE} تنها زمانی دادهٔ جدید نمونه‌برداری می‌شود که \code{new_data} فعال باشد. از آنجا که در حالت‌های \code{START} تا \code{STOP} چنین شرطی وجود ندارد، درخواست مجدد \code{new_data} در زمان \code{busy} نمی‌تواند قاب جاری را تغییر دهد. این تصمیم طراحی در آزمون هشتم به‌طور مستقیم ارزیابی شده است.

\begin{figure}[H]
  \centering
  \begin{tikzpicture}[
    >=Latex,
    statebox/.style={draw, rounded corners=3pt, line width=0.9pt, align=center,
      minimum width=2.65cm, minimum height=1.55cm, inner sep=4pt, fill=black!2},
    trans/.style={-Latex, line width=0.95pt},
    lbl/.style={font=\scriptsize, align=center, fill=white, inner sep=1.5pt},
    note/.style={draw, dashed, rounded corners=2pt, align=center, font=\scriptsize,
      inner sep=4pt, fill=black!1}
  ]
    \node[statebox] (idle)   at (0,0)   {\lr{\textbf{IDLE}}\\[-1pt]\scriptsize \code{tx=1}\\\code{busy=0}};
    \node[statebox] (start)  at (4.60,0) {\lr{\textbf{START}}\\[-1pt]\scriptsize \code{tx=0}\\\code{busy=1}};
    \node[statebox] (parity) at (9.20,0) {\lr{\textbf{PARITY}}\\[-1pt]\scriptsize \code{tx=parity_reg}\\\code{busy=1}};
    \node[statebox,minimum width=3.30cm] (data) at (9.20,-3.60) {\lr{\textbf{DATA}}\\[-1pt]\scriptsize \code{tx=data_reg[bit_idx]}\\\code{busy=1}};
    \node[statebox] (stop)   at (3.60,-3.60) {\lr{\textbf{STOP}}\\[-1pt]\scriptsize \code{tx=1}\\\code{busy=1}};

    \draw[trans] (-1.85,0) -- (idle.west);
    \node[lbl] at (-1.45,0.32) {\rl{شروع}};

    \draw[trans] (idle.east) -- node[lbl,above]{\code{new_data=1}} (start.west);
    \draw[trans] (start.east) -- node[lbl,above]{\rl{پایان زمان بیت}} (parity.west);
    \draw[trans] (parity.south) -- node[lbl,right]{\rl{پایان زمان بیت}} (data.north);
    \draw[trans] (data.west) -- node[lbl,above,text width=2.25cm]{\rl{پایان زمان بیت}\\\code{bit_idx=6}} (stop.east);
    \draw[trans] (stop.west) to[out=180,in=-90]
      node[lbl,pos=0.50,above left]{\rl{پایان زمان بیت}} (idle.south);

    \draw[trans] (data.south east) to[out=-25,in=25,looseness=2.45]
      node[lbl,right,text width=2.20cm]{\rl{پایان زمان بیت}\\\code{bit_idx<6}} (data.north east);

    \node[note,text width=6.3cm] at (0.90,-5.65) {\rl{با \code{rstN=0}: بازنشانی ناهم‌زمان و بازگشت به \lr{IDLE}}\\\rl{پایان زمان بیت: \code{tick_count=BIT_TICKS-1}}};
  \end{tikzpicture}
  \caption{نمودار ماشین حالت فرستندۀ \lr{UART}}
  \label{fig:sender-fsm}
\end{figure}

\subsection{کد \lr{Verilog} فرستنده}

کد کامل پیمانۀ فرستنده در فایل \filename{UARTSender.v} قرار دارد:

\begin{latin}
\lstinputlisting{../quartus/UARTSender.v}
\end{latin}

\begin{center}
\captionof{lstlisting}{پیاده‌سازی فرستندۀ \lr{UART}}
\label{lst:sender}
\end{center}

در کد، چهار جزء اصلی مسئول تولید قاب هستند: ثبات \code{data_reg} برای نگهداری دادهٔ لچ‌شده، ثبات \code{parity_reg} برای بیت توازن، شمارندۀ \code{tick_count} برای زمان‌بندی و \code{bit_idx} برای انتخاب بیت جاری. خروجی \code{tx} در هر حالت مستقیماً از وضعیت قاب تعیین می‌شود و \code{busy} نیز در تمام حالات غیر از \code{IDLE} فعال باقی می‌ماند. بنابراین ماشین حالت علاوه بر کنترل ترتیب منطقی قاب، کنترل زمانی انتقال را نیز بر عهده دارد.

\section{طراحی گیرنده}

\subsection{ساختار ماشین حالت گیرنده}

گیرنده نیز با یک ماشین حالت متناهی پنج‌حالته پیاده‌سازی شده است که همان ساختار کلی قاب فرستنده را دنبال می‌کند؛ با این تفاوت که وظیفۀ آن نمونه‌برداری از خط ورودی و بازسازی تدریجی قاب است. حالت‌های اصلی عبارت‌اند از \code{IDLE}، \code{START}، \code{PARITY}، \code{DATA} و \code{STOP}. علاوه بر ثبات حالت، گیرنده از ثبات \code{data_reg} برای ذخیرۀ دادهٔ دریافتی، \code{parity_reg} برای نگهداری parity دریافت‌شده، \code{bit_idx} برای مشخص‌کردن شمارهٔ بیت و \code{tick_count} برای کنترل زمان نمونه‌برداری استفاده می‌کند.

در حالت \code{IDLE} خط \code{rx} پیوسته بررسی می‌شود و تا زمانی که مقدار آن ۱ باشد، گیرنده در همین حالت باقی می‌ماند. مشاهدهٔ یک صفر نخستین نشانهٔ شروع قاب است و باعث می‌شود ماشین وارد حالت \code{START} شود. نکتهٔ مهم این است که این صفر به‌تنهایی به‌عنوان start معتبر پذیرفته نمی‌شود. گیرنده باید پس از گذشت نیم زمان بیت دوباره خط را بررسی کند تا مشخص شود سطح صفر واقعاً مربوط به بیت شروع بوده است. برای این منظور پارامتر محلی \code{HALF_BIT_TICKS} از \code{BIT_TICKS} به‌دست می‌آید که در پیاده‌سازی حاضر برابر دو چرخه است.

این تأیید میانی نقش مهمی در حذف پالس‌های کاذب دارد. اگر ورودی فقط برای مدت کوتاهی صفر شود و پیش از رسیدن به مرکز بیت دوباره یک شود، شرط حالت \code{START} برقرار نمی‌شود و ماشین مستقیماً به \code{IDLE} بازمی‌گردد. در نتیجه هیچ parity، داده یا stop bit اشتباهی دریافت نمی‌شود. همین رفتار در آزمون دهم، یعنی آزمون false start، مستقیماً بررسی شده است.

پس از تأیید start، گیرنده وارد حالت \code{PARITY} می‌شود. در این مرحله یک زمان بیت کامل صبر می‌کند و سپس مقدار \code{rx} را در \code{parity_reg} ذخیره می‌کند. پس از ذخیرۀ parity، اندیس \code{bit_idx} از صفر شروع می‌شود تا دریافت هفت بیت داده را کنترل کند.

در حالت \code{DATA}، در پایان هر زمان بیت مقدار فعلی \code{rx} در خانهٔ متناظر \code{data_reg[bit_idx]} نوشته می‌شود. بنابراین اولین نمونه در \code{data_reg[0]} و آخرین نمونه در \code{data_reg[6]} قرار می‌گیرد. این ترتیب دقیقاً با ترتیب ارسال فرستنده سازگار است و باعث می‌شود پس از دریافت آخرین بیت، \code{data_reg} همان مقدار هفت‌بیتی ارسال‌شده را در خود داشته باشد. اگر اندیس کمتر از ۶ باشد، تنها اندیس افزایش پیدا می‌کند و دریافت ادامه می‌یابد. پس از دریافت \code{D6}، ماشین به حالت \code{STOP} منتقل می‌شود.

در حالت \code{STOP} گیرنده پس از پایان زمان بیت، ابتدا مقدار خط را به‌عنوان stop bit بررسی می‌کند. اگر \code{rx=1} باشد، شرط نخست اعتبار قاب برقرار است. هم‌زمان parity محاسبه‌شده از هفت بیت داده با مقدار ذخیره‌شده در \code{parity_reg} مقایسه می‌شود. شرط نهایی در واقع ترکیب این دو بررسی است:

\[
(rx=1)\quad\land\quad(\oplus data\_reg=parity\_reg)
\]

تنها در صورت برقرار بودن هر دو شرط، \code{correct_data} برابر یک می‌شود؛ در غیر این صورت مقدار آن صفر خواهد بود. بنابراین قاب دارای parity اشتباه یا stop bit صفر، حتی اگر دادهٔ خام آن دریافت شده باشد، به‌عنوان قاب صحیح گزارش نمی‌شود.

صرف‌نظر از معتبر یا نامعتبر بودن قاب، دادهٔ موجود در \code{data_reg} در پایان حالت \code{STOP} به \code{rec_data} منتقل می‌شود و \code{rec_new_data} برای یک چرخه فعال می‌گردد. این جداسازی مهم است: \code{rec_new_data} اعلام می‌کند که فرآیند دریافت یک قاب تمام شده است، درحالی‌که \code{correct_data} مشخص می‌کند آیا آن قاب از نظر شرایط صحت معتبر بوده یا خیر. بنابراین مصرف‌کنندهٔ مدار می‌تواند بین «داده دریافت شده» و «داده معتبر» تمایز بگذارد.

پس از پایان \code{STOP}، گیرنده دوباره به \code{IDLE} بازمی‌گردد و برای قاب بعدی آماده می‌شود. همچنین در ابتدای هر چرخهٔ عادی، \code{rec_new_data} به‌صورت پیش‌فرض صفر می‌شود؛ به همین دلیل این سیگنال ذاتاً یک پالس یک‌کلاکی است و در آزمون‌های فایل آزمون نیز این ویژگی به‌صورت مستقل بررسی شده است.

بازنشانی گیرنده نیز فعال‌پایین و ناهم‌زمان است. با اعمال \code{rstN=0}، حالت، داده، parity، اندیس بیت و شمارندۀ زمان همگی به مقادیر اولیه بازمی‌گردند و \code{rec_new_data} و \code{correct_data} فوراً صفر می‌شوند. در نتیجه اگر بازنشانی وسط دریافت اتفاق بیفتد، قاب ناقص نمی‌تواند به‌عنوان یک دریافت کامل یا صحیح ثبت شود. این رفتار در آزمون نهم به‌صورت مستقیم ارزیابی شده است.

\begin{figure}[H]
  \centering
  \begin{tikzpicture}[
    >=Latex,
    statebox/.style={draw, rounded corners=3pt, line width=0.9pt, align=center,
      minimum width=2.75cm, minimum height=1.55cm, inner sep=4pt, fill=black!2},
    trans/.style={-Latex, line width=0.95pt},
    lbl/.style={font=\scriptsize, align=center, fill=white, inner sep=1.5pt},
    note/.style={draw, dashed, rounded corners=2pt, align=center, font=\scriptsize,
      inner sep=4pt, fill=black!1}
  ]
    \node[statebox] (idle)   at (0,0)   {\lr{\textbf{IDLE}}\\[-1pt]\scriptsize \rl{انتظار برای \code{rx=0}}};
    \node[statebox,minimum width=3.05cm] (start)  at (4.60,0) {\lr{\textbf{START}}\\[-1pt]\scriptsize \rl{تأیید شروع}\\\rl{در نیمهٔ بیت}};
    \node[statebox,minimum width=3.05cm] (parity) at (9.20,0) {\lr{\textbf{PARITY}}\\[-1pt]\scriptsize \rl{نمونه‌برداری و ذخیرهٔ}\\\rl{بیت توازن}};
    \node[statebox,minimum width=3.20cm] (data)   at (9.20,-3.60) {\lr{\textbf{DATA}}\\[-1pt]\scriptsize \rl{نمونه‌برداری از \code{D0} تا \code{D6}}};
    \node[statebox,minimum width=3.10cm] (stop)   at (3.60,-3.60) {\lr{\textbf{STOP}}\\[-1pt]\scriptsize \rl{بررسی بیت توقف}\\\rl{و توازن}};

    \draw[trans] (-1.85,0) -- (idle.west);
    \node[lbl] at (-1.45,0.32) {\rl{شروع}};

    \draw[trans] (idle.east) -- node[lbl,above]{\code{rx=0}} (start.west);
    \draw[trans] (start.east) -- node[lbl,above,text width=2.05cm]{\rl{بررسی نیم‌بیت}\\\code{rx=0}} (parity.west);
    \draw[trans] (start.north west) to[out=145,in=35,looseness=1.20]
      node[lbl,above,text width=3.1cm]{\rl{بررسی نیم‌بیت، \code{rx=1}}\\\rl{شروع کاذب}} (idle.north east);
    \draw[trans] (parity.south) -- node[lbl,right]{\rl{پایان زمان بیت}} (data.north);
    \draw[trans] (data.west) -- node[lbl,above,text width=2.25cm]{\rl{پایان زمان بیت}\\\code{bit_idx=6}} (stop.east);
    \draw[trans] (stop.west) to[out=180,in=-90]
      node[lbl,pos=0.50,above left]{\rl{پایان زمان بیت}} (idle.south);

    \draw[trans] (data.south east) to[out=-25,in=25,looseness=2.45]
      node[lbl,right,text width=2.20cm]{\rl{پایان زمان بیت}\\\code{bit_idx<6}} (data.north east);

    \node[note,text width=6.2cm] at (0.70,-5.85) {\rl{با \code{rstN=0}: بازنشانی ناهم‌زمان و بازگشت به \lr{IDLE}}\\\rl{پایان زمان بیت: \code{tick_count=BIT_TICKS-1}}\\\rl{بررسی نیم‌بیت: \code{tick_count=HALF_BIT_TICKS-1}}};
    \node[note,text width=5.25cm] at (8.65,-5.85) {\rl{در پایان \lr{STOP}: قاب فقط در صورت \code{rx=1} و تطابق توازن معتبر است.}\\\code{rec_new_data=1}};
  \end{tikzpicture}
  \caption{نمودار ماشین حالت گیرندۀ \lr{UART}}
  \label{fig:receiver-fsm}
\end{figure}

\subsection{کد \lr{Verilog} گیرنده}

کد کامل پیمانۀ گیرنده در فایل \filename{UARTReceiver.v} آمده است:

\begin{latin}
\lstinputlisting{../quartus/UARTReceiver.v}
\end{latin}

\begin{center}
\captionof{lstlisting}{پیاده‌سازی گیرندۀ \lr{UART}}
\label{lst:receiver}
\end{center}

در این کد، \code{HALF_BIT_TICKS} برای تأیید مرکز start، \code{tick_count} برای نگهداری زمان هر بیت، \code{bit_idx} برای ترتیب ذخیرۀ هفت بیت و \code{parity_reg} برای نگهداری parity دریافتی به‌کار می‌روند. در نتیجه کل فرآیند دریافت، از تشخیص شروع تا بررسی نهایی قاب، تنها با یک ماشین حالت متناهی و مجموعه‌ای کوچک از ثبات‌های زمان‌بندی و ذخیره‌سازی کنترل می‌شود.

\section{پیمانهٔ سطح بالا و اتصال بازگشت مستقیم}

برای آزمایش مسیر کامل انتقال، دو پیمانهٔ فرستنده و گیرنده در یک پیمانهٔ سطح بالا به نام \filename{UARTTop.v} نمونه‌سازی شده‌اند. خروجی سریال فرستنده در سیم داخلی \code{rx_line} قرار می‌گیرد و همان سیم بدون هیچ عنصر میانجی به ورودی \code{rx} گیرنده متصل می‌شود. بنابراین ساختار بازگشت مستقیم ایجاد شده است.

این اتصال دو مزیت مهم دارد: نخست آنکه نیازی به تولید یک مسیر سیمی خارجی برای آزمایش دریافت وجود ندارد و دوم آنکه کل زنجیرۀ ارسال و دریافت به‌صورت یک واحد واحد قابل آزمون است. در حالت صحیح، هر داده‌ای که فرستنده ارسال کند باید در خروجی گیرنده همان داده را تولید کند.

\subsection{کد \lr{Verilog} پیمانۀ سطح بالا}

کد کامل پیمانۀ سطح بالا:

\begin{latin}
\lstinputlisting{../quartus/UARTTop.v}
\end{latin}

\begin{center}
\captionof{lstlisting}{اتصال فرستنده و گیرنده در پیمانۀ \lr{UARTTop}}
\label{lst:top}
\end{center}

\section{سنتز و مشاهدۀ ساختار در \lr{Quartus}}

\subsection{روند انجام کار در \lr{Quartus}}

برای مشاهدهٔ ساختار سخت‌افزاری حاصل از کدهای \lr{Verilog}، ابتدا پروژه در \lr{Quartus} ایجاد شد و فایل‌های \filename{UARTSender.v}، \filename{UARTReceiver.v} و \filename{UARTTop.v} به پروژه افزوده شدند. سپس پیمانۀ \code{UARTTop} به‌عنوان موجودیت سطح بالا انتخاب شد. پس از آن تحلیل و سنتز پروژه انجام گرفت تا قابل‌سنتزبودن طراحی بررسی شود.

مسیر مورد استفاده برای بازکردن نمای سطح انتقال ثبات\enfoot{Register-Transfer Level (RTL)} به‌صورت زیر است:

\begin{latin}
\begin{center}
\texttt{Tools > Netlist Viewers > RTL Viewer}
\end{center}
\end{latin}

پس از اجرای این گزینه، \lr{Quartus} ساختار منطقی استخراج‌شده از توصیف \lr{Verilog} را نمایش می‌دهد. در این نما، ثبات‌ها، منطق انتخاب، شمارنده‌ها، مسیرهای داده و روابط بین آن‌ها قابل مشاهده هستند. در گزارش حاضر سه تصویر مستقل از ساختارهای تولیدشده برای پیمانه‌های اصلی آورده شده است.

\subsection{نمای \lr{RTL} فرستنده}

\begin{figure}[H]
  \centering
  \reportimage{images/image01.png}
  \caption{نمای \lr{RTL} حاصل از سنتز پیمانۀ \lr{UARTSender}}
  \label{fig:sender-rtl}
\end{figure}

در شکل \ref{fig:sender-rtl}، ساختار منطقی حاصل از توصیف رفتاری فرستنده دیده می‌شود. شلوغی این نما طبیعی است؛ زیرا ماشین حالت، شمارندۀ زمان‌بندی، ثبات داده، محاسبۀ توازن و منطق انتخاب بیت جاری پس از سنتز به مجموعه‌ای از عناصر منطقی صریح تبدیل شده‌اند. وجود مسیرهای متعدد در بخش پایینی شکل با تصمیم طراحی برای نگهداری حالت، شمارش چرخه‌ها و انتخاب \code{data_reg[bit_idx]} مرتبط است.

\subsection{نمای \lr{RTL} گیرنده}

\begin{figure}[H]
  \centering
  \reportimage{images/image02.png}
  \caption{نمای \lr{RTL} حاصل از سنتز پیمانۀ \lr{UARTReceiver}}
  \label{fig:receiver-rtl}
\end{figure}

شکل \ref{fig:receiver-rtl} ساختار گیرنده را نشان می‌دهد. در این نما علاوه بر ثبات‌های مربوط به حالت و شمارندۀ زمان، مسیر مربوط به ذخیرۀ داده و بیت توازن و همچنین منطق بررسی شرط نهایی قاب دیده می‌شود. پیچیدگی بیشتر ساختار گیرنده نسبت به یک منطق ترکیبی ساده، ناشی از دریافت مرحله‌به‌مرحلۀ قاب و کنترل زمان نمونه‌برداری است.

\subsection{نمای \lr{RTL} پیمانۀ سطح بالا}

\begin{figure}[H]
  \centering
  \reportimage{images/image03.png}
  \caption{نمای \lr{RTL} پیمانۀ سطح بالا و اتصال مستقیم فرستنده به گیرنده}
  \label{fig:top-rtl}
\end{figure}

شکل \ref{fig:top-rtl} نمایی بسیار خواناتر از ساختار سطح بالا ارائه می‌کند. در سمت چپ ورودی‌های ساعت، بازنشانی، درخواست داده و دادهٔ ارسالی وارد فرستنده می‌شوند. خروجی سریال فرستنده به‌طور مستقیم به ورودی سریال گیرنده وصل شده است و در سمت راست نیز \code{rec_data}، \code{rec_new_data} و \code{correct_data} مشاهده می‌شوند. این شکل دقیقاً رابط اصلی‌ای را که در شبیه‌سازی آزموده شده است نمایش می‌دهد.

\section{طراحی فایل آزمون}

\subsection{هدف و ساختار فایل \lr{Testbench}}

فایل آزمون\enfoot{Testbench} در \filename{Tester.v} طراحی شده و تنها به یک سناریوی سادهٔ ارسال و دریافت محدود نیست. در این فایل دو پیمانۀ تحت آزمون وجود دارد: نخست \code{top_dut} که کل مسیر فرستنده، بازگشت مستقیم و گیرنده را شامل می‌شود و دوم یک گیرندۀ مستقل به نام \code{manual_rx_dut}. گیرندۀ دوم عمداً به یک ورودی مستقل وصل شده است تا بتوان قاب‌های خراب را مستقیماً به آن اعمال کرد. به این ترتیب خطای توازن و خطای بیت توقف، بدون وابستگی به درست یا غلط بودن فرستنده، قابل آزمون هستند.

\subsection{تولید ساعت، بازنشانی و ثبت شکل موج}

در فایل آزمون، پارامتر \code{BIT_TICKS=4} و دورۀ ساعت \code{CLK_PERIOD=10} نانوثانیه تعیین شده‌اند. ساعت به‌صورت پیوسته با یک فرایند \code{forever} تولید می‌شود. بازنشانی فعال‌پایین و ناهم‌زمان است و تابع \code{reset_all} بلافاصله پس از اعمال صفر به \code{rstN}، وضعیت خروجی‌های مهم را بررسی می‌کند.

برای تولید فایل قابل مشاهده در نرم‌افزار تحلیل شکل موج، خروجی \lr{VCD}\enfoot{Value Change Dump} با نام \filename{uart_tests.vcd} ساخته می‌شود. در این فایل فقط سیگنال‌های لازم برای گزارش ذخیره شده‌اند تا نمای حاصل خوانا باقی بماند. ازجمله سیگنال‌های ذخیره‌شده می‌توان به \code{clk}، \code{rstN}، \code{new_data}، \code{send_data}، \code{tx}، \code{busy}، \code{rec_data}، \code{rec_new_data} و \code{correct_data} اشاره کرد. همچنین مقدار \code{test_id} عمداً در فایل شکل موج ذخیره شده است تا محدودهٔ هر آزمون در تحلیل تصویری به‌وضوح قابل تشخیص باشد.

\subsection{وظایف کمکی و خودآزمایی}

برای جلوگیری از تکرار منطق آزمون، چند وظیفۀ کمکی در فایل تعریف شده است. مهم‌ترین آن‌ها عبارت‌اند از:

\begin{itemize}[itemsep=0.3em]
  \item \code{add_check}: یک شرط را ثبت می‌کند، شمارندۀ بررسی‌ها را افزایش می‌دهد و در صورت شکست، شمارندۀ خطا را نیز افزایش می‌دهد.
  \item \code{reset_all}: بازنشانی اولیه را اعمال کرده و وضعیت بی‌کار فرستنده و گیرنده را کنترل می‌کند.
  \item \code{send_top_data}: پس از اطمینان از غیرمشغول‌بودن فرستنده، داده را در ورودی قرار داده و یک پالس کوتاه \code{new_data} تولید می‌کند.
  \item \code{check_top_data}: منتظر پالس \code{rec_new_data} می‌ماند و سپس هم داده و هم \code{correct_data} را با مقدار مورد انتظار مقایسه می‌کند. یک‌کلاک‌بودن این پالس نیز بررسی می‌شود.
  \item \code{send_manual_frame}: یک قاب کامل را به‌صورت دستی روی ورودی گیرندۀ مستقل اعمال می‌کند تا بتوان مقدار parity و stop را عمداً تغییر داد.
  \item \code{check_serial_frame}: تک‌تک بیت‌های روی \code{tx} و همچنین فعال‌بودن \code{busy} را در تمام چرخه‌های هر بیت بررسی می‌کند؛ بنابراین یک تست صرفاً مبتنی بر loopback نمی‌تواند خطای مشترک فرستنده و گیرنده را پنهان کند.
\end{itemize}

\section{سناریوهای آزمون و حالت‌های مرزی}

مجموعهٔ آزمون شامل ۱۱ سناریوی اصلی است. این سناریوها فقط مسیر عادی ارسال و دریافت را بررسی نمی‌کنند، بلکه رفتار مدار را در برابر خطا، مرزهای داده، درخواست‌های نامعتبر و رخدادهای غیرعادی نیز ارزیابی می‌کنند. در ادامه هر یک از آزمون‌ها به‌ترتیب اجرای فایل \filename{Tester.v} توضیح داده می‌شود.

\subsection{آزمون ۱: بازنشانی و حالت بی‌کار}

در نخستین آزمون، پس از اعمال بازنشانی، وضعیت اولیهٔ کل زنجیرۀ انتقال و دریافت بررسی می‌شود. در حالت بازنشانی انتظار می‌رود فرستنده بلافاصله به حالت بی‌کار برگردد؛ بنابراین \code{tx} باید برابر ۱ و \code{busy} برابر صفر باشد. هم‌زمان، سیگنال‌های وضعیت گیرنده یعنی \code{rec_new_data} و \code{correct_data} نیز باید صفر باشند. این آزمون علاوه بر بررسی وضعیت بی‌کار، رفتار بازنشانی ناهم‌زمان را نیز کنترل می‌کند؛ یعنی مدار نباید برای مشاهدهٔ لبهٔ ساعت صبر کند تا اثر \code{rstN=0} اعمال شود.

\subsection{آزمون ۲: دریافت صحیح، قالب قاب و زمان‌بندی سیگنال busy}

در این آزمون دادهٔ هفت‌بیتی \code{1011001} از مسیر اصلی \code{UARTTop} ارسال می‌شود. فایل آزمون سه بررسی را به‌صورت هم‌زمان انجام می‌دهد: نخست داده به فرستنده داده می‌شود، دوم قالب کامل قاب مستقیماً روی \code{tx} کنترل می‌شود و سوم گیرنده منتظر \code{rec_new_data} می‌ماند تا دادهٔ نهایی را با مقدار مورد انتظار مقایسه کند. بنابراین آزمون فقط برابر بودن دادهٔ ارسال‌شده و دریافت‌شده را بررسی نمی‌کند، بلکه ترتیب بیت‌ها، بیت شروع، parity، هفت بیت داده، بیت توقف و فعال‌بودن \code{busy} در تمام مدت قاب نیز کنترل می‌شوند. در پایان، \code{busy} باید تنها پس از پایان بیت توقف صفر شود و \code{tx} دوباره به وضعیت بی‌کار یعنی ۱ بازگردد.

\subsection{آزمون ۳: خطای توازن}

برای آزمون خطای parity از گیرندۀ مستقل \code{manual_rx_dut} استفاده شده است. داده همچنان \code{1011001} است، اما بیت توازن عمداً معکوس می‌شود. به این ترتیب دادهٔ دریافت‌شده می‌تواند در ثبات گیرنده ثبت شود، ولی شرط اعتبار قاب دیگر برقرار نیست. انتظار می‌رود \code{manual_correct_data} صفر شود. این ساختار آزمون باعث می‌شود تشخیص خطای parity مستقیماً در خود گیرنده سنجیده شود و نتیجه به عملکرد فرستندهٔ اصلی وابسته نباشد.

\subsection{آزمون ۴: خطای بیت توقف}

در آزمون چهارم، parity مطابق مقدار صحیح داده تولید می‌شود، اما بیت توقف به‌طور عمدی صفر قرار داده می‌شود. بنابراین حتی در صورت درست بودن داده و parity، قاب نباید معتبر اعلام شود. گیرنده در پایان قاب داده را ثبت می‌کند و رویداد دریافت را صادر می‌کند، ولی \code{manual_correct_data} باید برابر صفر باقی بماند. این آزمون نشان می‌دهد که شرط اعتبار قاب در پیاده‌سازی فقط به parity محدود نیست و مقدار بیت توقف نیز به‌طور مستقل بررسی می‌شود.

\subsection{آزمون ۵: ارسال چند قاب پشت‌سرهم}

در این سناریو دو دادهٔ متفاوت یعنی \code{1011001} و \code{0101101} در دو قاب جداگانه ارسال می‌شوند. فایل آزمون پس از پایان قاب اول منتظر غیرفعال شدن \code{busy} می‌ماند و سپس قاب دوم را آغاز می‌کند. در سمت دریافت نیز دو بار منتظر پالس \code{rec_new_data} می‌ماند و هر دو مقدار را جداگانه با دادهٔ مورد انتظار مقایسه می‌کند. هدف این آزمون اطمینان از آن است که پایان یک قاب و آغاز قاب بعدی به‌درستی از یکدیگر تفکیک شوند و یک رویداد دریافت به‌جای دو قاب ایجاد نشود.

\subsection{آزمون ۶: حداقل داده}

در این آزمون کمینهٔ عددی فضای ورودی یعنی \code{0000000} ارسال می‌شود. برای این داده parity برابر صفر است، بنابراین علاوه بر بررسی مسیر معمول دریافت، حالت مرزی parity=0 نیز آزمایش می‌شود. انتظار می‌رود \code{rec_data} برابر \code{0000000} و \code{correct_data} برابر یک شود. این سناریو ثابت می‌کند که مدار در مرز پایین فضای داده نیز بدون ایجاد حالت خاص یا خطای محاسباتی عمل می‌کند.

\subsection{آزمون ۷: حداکثر داده}

در آزمون هفتم بیشینهٔ مقدار ممکن برای دادهٔ هفت‌بیتی یعنی \code{1111111} ارسال می‌شود. در این حالت parity برابر یک است. این آزمون مکمل آزمون ششم است و علاوه بر بررسی مرز بالای داده، مسیر معتبر parity=1 را نیز فعال می‌کند؛ در نتیجه دو مقدار مرزی و هر دو حالت اصلی parity در آزمون‌های جهت‌دار پوشش داده می‌شوند.

\subsection{آزمون ۸: فعال‌سازی \lr{new\_data} هنگام مشغول بودن فرستنده}

در آزمون هشتم ابتدا دادهٔ \code{0010110} برای آغاز یک قاب در فرستنده قرار می‌گیرد. پس از آن، در حالی که \code{busy=1} است، مقدار ورودی به \code{1101001} تغییر داده شده و یک پالس جدید \code{new_data} اعمال می‌شود. رفتار مورد انتظار آن است که درخواست جدید نادیده گرفته شود و دادهٔ ذخیره‌شدهٔ قاب جاری تغییر نکند. فایل آزمون علاوه بر بررسی دریافت صحیح قاب اول، تعداد رخدادهای دریافت را نیز می‌شمارد تا اطمینان حاصل شود که درخواست زمان مشغول بودن، قاب دوم اضافی ایجاد نکرده است. بنابراین این آزمون هم لچ شدن داده در آغاز انتقال و هم حفاظت از قاب جاری را بررسی می‌کند.

\subsection{آزمون ۹: بازنشانی ناهم‌زمان در میانهٔ انتقال و بازیابی}

در آزمون نهم یک قاب با دادهٔ \code{0110101} شروع می‌شود و پس از چند چرخه، بازنشانی ناهم‌زمان در وسط قاب فعال می‌گردد. انتظار می‌رود فرستنده فوراً انتقال را قطع کرده و به وضعیت \code{tx=1, busy=0} برگردد. در گیرنده نیز وضعیت‌های دریافت باید پاک شوند و قاب ناقص نباید به‌عنوان یک رویداد دریافت معتبر ثبت شود. پس از آزادشدن بازنشانی، یک قاب جدید با دادهٔ \code{1000011} ارسال می‌شود تا توانایی مدار در بازیابی کامل پس از قطع ناگهانی نیز اثبات شود.

\subsection{آزمون ۱۰: رد شروع کاذب و بازیابی گیرنده}

در آزمون دهم ابتدا یک پالس کوتاه صفر روی \code{manual_rx} ایجاد می‌شود که طول آن از زمان لازم برای تأیید بیت شروع کمتر است. گیرنده باید این رخداد را یک شروع واقعی در نظر نگیرد و هیچ رویداد دریافت جدیدی تولید نکند. سپس یک قاب معتبر با دادهٔ \code{0000001} و parity صحیح ارسال می‌شود تا مشخص شود رد شدن شروع کاذب باعث قفل‌شدن گیرنده نشده و گیرنده می‌تواند بلافاصله به حالت عادی بازگردد. این آزمون مستقیماً منطق تأیید start در نزدیکی مرکز بیت را ارزیابی می‌کند.

\subsection{آزمون ۱۱: پیمایش جامع تمام ۱۲۸ مقدار داده}

در آخرین آزمون، تمام مقادیر ممکن یک دادهٔ هفت‌بیتی از صفر تا ۱۲۷ به‌صورت پشت‌سرهم ارسال می‌شوند. این آزمون از نظر پوشش فضای داده کامل است و به‌جای تکیه بر چند نمونهٔ خاص، رفتار فرستنده و گیرنده را برای تمام ترکیب‌های ممکن داده بررسی می‌کند. برای جلوگیری از شلوغ شدن فایل شکل موج، ثبت \lr{VCD} در طول این آزمون موقتاً غیرفعال می‌شود، اما بررسی‌های منطقی همچنان انجام می‌شوند. در پایان باید تمام ۱۲۸ مقدار بدون خطا دریافت شوند.

آزمون‌های بالا در مجموع مسیر صحیح، دو نوع خطای قاب، مرزهای داده، درخواست در زمان مشغول بودن، بازنشانی میانهٔ قاب، شروع کاذب و کل فضای دادهٔ ورودی را پوشش می‌دهند. بنابراین اعتبار نتیجهٔ نهایی تنها بر یک شکل موج یا یک دادهٔ نمونه متکی نیست.

\section{شکل موج‌های منتخب و تحلیل آن‌ها}

از آنجا که اجرای کامل فایل آزمون شامل ۱۱ سناریو و در آزمون آخر ۱۲۸ قاب است، نمایش همۀ بازه‌های زمانی در یک شکل واحد باعث شلوغی شدید نمودار می‌شود. به همین دلیل در گزارش، شکل موج‌های منتخب مربوط به آزمون‌های مهم آورده شده‌اند. با این حال، پیش از نمایش این تصاویر، همۀ \code{test_id}ها و هدف آن‌ها در بخش قبل توضیح داده شده‌اند؛ بنابراین شکل‌های حاضر نمونه‌هایی از اجرای همان سناریوها هستند و جایگزین آزمون‌های اجراشده نیستند.

\subsection{شکل موج آزمون دریافت صحیح، \lr{test\_id=2}}

در این شکل آغاز انتقال با پالس \code{new_data} مشخص می‌شود و پس از آن \code{busy} فعال شده و \code{tx} به‌ترتیب بیت‌های قاب را حمل می‌کند. مقدار \code{send_data=1011001} در فرستنده لچ شده است و در پایان دریافت، \code{rec_data} همان مقدار را نشان می‌دهد. پالس کوتاه \code{rec_new_data} پایان فرآیند دریافت را مشخص می‌کند و فعال شدن \code{correct_data} نشان می‌دهد که قاب از نظر parity و stop bit معتبر بوده است.

\begin{figure}[H]
  \centering
  \reportimage[0.92\textwidth]{images/image04.png}
  \caption{شکل موج منتخب برای آزمون دریافت صحیح و بررسی قالب قاب، \code{test_id=2}}
  \label{fig:wave-test2}
\end{figure}

\subsection{شکل موج آزمون خطای parity، \lr{test\_id=3}}

در این بازه قاب با دادهٔ صحیح ساخته می‌شود، اما parity به‌صورت عمدی نادرست است. گیرنده مستقل همچنان فرآیند دریافت را کامل می‌کند و داده را ذخیره می‌کند، ولی سیگنال صحت قاب فعال نمی‌شود. بنابراین شکل موج به‌طور هم‌زمان تفاوت میان «دریافت داده» و «معتبر بودن قاب» را نشان می‌دهد.

\begin{figure}[H]
  \centering
  \reportimage[0.92\textwidth]{images/image05.png}
  \caption{شکل موج منتخب برای آزمون خطای parity، \code{test_id=3}}
  \label{fig:wave-test3}
\end{figure}

\subsection{شکل موج آزمون خطای بیت توقف، \lr{test\_id=4}}

در این آزمون parity درست است، اما stop bit برابر صفر قرار داده شده است. به همین دلیل در پایان قاب، پالس دریافت ایجاد می‌شود ولی \code{manual_correct_data} صفر باقی می‌ماند. این شکل نشان می‌دهد که بررسی بیت توقف شرط مستقلی برای معتبر شناختن قاب است.

\begin{figure}[H]
  \centering
  \reportimage[0.92\textwidth]{images/image06.png}
  \caption{شکل موج منتخب برای آزمون خطای بیت توقف، \code{test_id=4}}
  \label{fig:wave-test4}
\end{figure}

\subsection{شکل موج آزمون ارسال چند قاب، \lr{test\_id=5}}

در این شکل دو قاب با داده‌های متفاوت پشت‌سرهم دیده می‌شوند. پس از پایان قاب نخست و بازگشت \code{busy} به صفر، قاب دوم آغاز می‌شود. در سمت گیرنده، هر داده به‌صورت مستقل در \code{rec_data} ظاهر شده و برای هر قاب یک پالس مجزا در \code{rec_new_data} ایجاد می‌شود. بنابراین زمان‌بندی پایان و آغاز قاب‌های متوالی صحیح است.

\begin{figure}[H]
  \centering
  \reportimage[0.92\textwidth]{images/image07.png}
  \caption{شکل موج منتخب برای آزمون ارسال چند قاب پشت‌سرهم، \code{test_id=5}}
  \label{fig:wave-test5}
\end{figure}

\subsection{شکل موج آزمون نادیده‌گرفتن \lr{new\_data} هنگام \lr{busy}، \lr{test\_id=8}}

در این شکل در حالی که فرستنده در حال ارسال قاب اول است، دادهٔ ورودی تغییر کرده و یک درخواست جدید اعمال می‌شود. با وجود این تغییر، قاب جاری با همان دادهٔ اولیه ادامه پیدا می‌کند و پس از پایان آن تنها یک رویداد دریافت ایجاد می‌شود. این رفتار تأیید می‌کند که داده در ابتدای انتقال لچ شده و درخواست‌های جدید هنگام \code{busy} باعث خراب شدن یا صف شدن قاب دیگری نشده‌اند.

\begin{figure}[H]
  \centering
  \reportimage[0.92\textwidth]{images/image08.png}
  \caption{شکل موج منتخب برای آزمون درخواست جدید هنگام مشغول‌بودن فرستنده، \code{test_id=8}}
  \label{fig:wave-test8}
\end{figure}

\subsection{شکل موج آزمون بازنشانی میانهٔ انتقال، \lr{test\_id=9}}

در این بازه ابتدا یک قاب در حال انتقال است، سپس \code{rstN} به صفر می‌رود. اثر بازنشانی ناهم‌زمان بلافاصله روی \code{tx} و \code{busy} دیده می‌شود و قاب جاری بدون تولید رویداد دریافت معتبر متوقف می‌شود. پس از آزادشدن بازنشانی، قاب جدید با دادهٔ \code{1000011} به‌صورت عادی منتقل شده و در گیرنده دریافت می‌شود؛ بنابراین مدار پس از یک قطع ناگهانی نیز به وضعیت صحیح برمی‌گردد.

\begin{figure}[H]
  \centering
  \reportimage[0.92\textwidth]{images/image09.png}
  \caption{شکل موج منتخب برای آزمون بازنشانی ناهم‌زمان و بازیابی، \code{test_id=9}}
  \label{fig:wave-test9}
\end{figure}

\subsection{توضیح کلی شکل موج‌ها}

در تمامی شکل‌های بالا، سیگنال \code{clk} مبنای زمان‌بندی مدار است و با توجه به مقدار \code{BIT_TICKS=4} هر بیت حدود ۴۰ نانوثانیه روی خط قرار می‌گیرد. سیگنال \code{tx} قاب را به‌ترتیب مورد انتظار حمل می‌کند، \code{busy} محدودهٔ زمانی مشغول بودن فرستنده را مشخص می‌سازد، و \code{rec_new_data} پایان دریافت هر قاب را به‌صورت یک پالس کوتاه اعلام می‌کند. مقایسۀ \code{send_data} و \code{rec_data} نیز امکان بررسی مستقیم صحت بازگشت داده را فراهم می‌کند.

تصاویر حاضر به‌صورت منتخب از میان بازه‌های اجرای آزمون تهیه شده‌اند؛ در نتیجه ممکن است آزمون‌های ۱، ۶، ۷، ۱۰ و ۱۱ در تصویر مستقل نداشته باشند، ولی منطق و نتیجهٔ همۀ آن‌ها در بخش سناریوها و خروجی نهایی فایل آزمون بررسی شده است.

\section{خروجی نهایی فایل آزمون}

آخرین تصویر، خروجی متنی فایل آزمون است و نتیجۀ خودآزمایی تمام سناریوها را نشان می‌دهد.

\begin{figure}[H]
  \centering
  \includegraphics[height=0.72\textheight,keepaspectratio]{images/image10.png}
  \caption{خروجی نهایی فایل آزمون و خلاصۀ اجرای آزمون‌های \lr{UART}}
  \label{fig:test-summary}
\end{figure}

مطابق این خروجی، همۀ آزمون‌های اصلی با پیام‌های \lr{PASS} پشت‌سر گذاشته شده‌اند. در آزمون آخر نیز تمام ۱۲۸ مقدار ممکن دادهٔ هفت‌بیتی بدون خطا دریافت شده‌اند. همچنین در خلاصۀ نهایی، تعداد کل بررسی‌ها برابر ۳۶۹ و تعداد خطاها برابر صفر گزارش شده است و در نتیجه عبارت \lr{ALL TESTS PASSED} نمایش داده می‌شود.

\section{تحلیل نتایج و جمع‌بندی}

نتایج این آزمایش سه بخش اصلی طراحی را تأیید می‌کند. نخست، منطق فرستنده و گیرنده مطابق قالب قاب تعریف‌شده عمل می‌کند؛ یعنی بیت شروع، parity، هفت بیت داده و بیت توقف با ترتیب و زمان‌بندی صحیح روی خط قرار می‌گیرند و گیرنده نیز همان ترتیب را دنبال می‌کند. دوم، اعتبار قاب به‌صورت مستقل از دریافت داده بررسی می‌شود و خطا در parity یا stop bit باعث صفرشدن سیگنال صحت می‌شود. سوم، فایل آزمون فقط رفتار عادی مدار را بررسی نمی‌کند و حالت‌های مرزی مؤثّر بر عملکرد واقعی مدار را نیز پوشش می‌دهد.

وجود گیرندۀ مستقل برای تزریق خطا باعث شده است آزمون‌های parity و stop bit واقعاً مستقل از منطق فرستنده باشند. بررسی مستقیم بیت‌های \code{tx} نیز از آن جلوگیری می‌کند که در یک بازگشت مستقیم، اگر فرستنده و گیرنده یک خطای مشترک در ترتیب بیت داشته باشند، نتیجه به‌اشتباه صحیح گزارش شود. همچنین بررسی یک‌کلاک‌بودن \code{rec_new_data} تضمین می‌کند که رویداد دریافت به‌صورت یک پالس قابل استفاده برای منطق بالاتر تولید می‌شود.

در نهایت، موفقیت همۀ ۳۶۹ بررسی و صفر بودن تعداد خطاها، همراه با عبور کامل از پیمایش ۱۲۸ مقدار دادهٔ هفت‌بیتی، نشان می‌دهد که طراحی ارائه‌شده برای سناریوهای بررسی‌شده پایدار و قابل‌اعتماد است. بنابراین می‌توان نتیجه گرفت که ساختار فرستنده، گیرنده، اتصال سطح بالا و فایل آزمون جامع، تمام خواسته‌های اصلی آزمایش را برآورده کرده‌اند.

\section{فایل‌های تحویلی}

مطابق صورت آزمایش، فایل‌های اصلی همراه گزارش عبارت‌اند از:

\begin{itemize}[itemsep=0.25em]
  \item \filename{UARTSender.v}: طراحی فرستنده و ماشین حالت ارسال؛
  \item \filename{UARTReceiver.v}: طراحی گیرنده، دریافت داده و بررسی parity و stop bit؛
  \item \filename{UARTTop.v}: اتصال فرستنده به گیرنده در قالب بازگشت مستقیم؛
  \item \filename{Tester.v}: فایل آزمون جامع برای تولید ساعت، اعمال بازنشانی، تولید داده، تزریق خطا، بررسی خروجی‌ها و اجرای آزمون جامع ۱۲۸ حالتی؛
  \item \filename{uart_tests.vcd}: فایل حاصل از ثبت تغییرات سیگنال‌ها در اجرای شبیه‌سازی.
\end{itemize}

\end{document}
